\subsection{Complete intersection ideals and regular rings}

PROPOSITION 2.--- Let A be a Noetherian ring. The following conditions are equivalent:

\begin{enumerate}
\def\labelenumi{(\roman{enumi})}
\item
  A is regular;
\item
  every maximal ideal of A is completely secant;
\item
  Every ideal J of A such that A/J is regular is completely secant.
\end{enumerate}

(i)⇒(iii): suppose the ring A is regular; let J be an ideal of A such that Λ/J is regular, and let p be a prime ideal of A containing J. Then the local rings \(A_{p}\) and \(A_{p}/J_{p}\) are regular, hence \(J_{p}\) is generated by a completely secant sequence for \(A_{p}\) (VIII, § 5, no. 3, cor. 2 and prop. 2), which means that J is completely secant at p.

\begin{enumerate}
\def\labelenumi{(\roman{enumi})}
\setcounter{enumi}{2}
\tightlist
\item
  (ii): This is clear since a field is a regular ring.
\end{enumerate}

(ii)⇒(i): let m be a maximal ideal of A; under hypothesis (ii), the maximal ideal mA \(_{m}\) of Λ \(_{m}\) is generated by a completely secant sequence for A \(_{m}\) therefore Λ \(_{m}\) is regular (VIII, § 5, no. 2, th. 1). Consequently, A is regular.

PROPOSITION 3.--- Let A be a Noetherian ring and J an ideal of A such that A/J is regular. The following conditions are equivalent:

\begin{enumerate}
\def\labelenumi{(\roman{enumi})}
\item
  the ideal J is completely secant ;
\item
  for every prime (resp. maximal) ideal p of A containing J, the ring A \(_{p}\) is regular;
\item
  The separated completion of A with respect to the J-adic topology is regular.
\end{enumerate}

By th. 1, condition (i) means that for every prime (resp. maximal) ideal p of A containing J, the ideal \(J_{p}\) of the local ring \(A_{p}\) is generated by a completely secant sequence for \(A_{p}\) . Since \(A_{p}/J_{p}\) is regular by hypothesis; this latter condition is equivalent to the fact that \(A_{p}\) let it be regular (VIII, § 5, no. 3, prop. 2 and its cor. 2); this proves the equivalence of (i) and (ii). The equivalence of (ii) and (iii) follows from § 4, no. 2, cor. 3 of prop. 4.

PROPOSITION 4.--- Let A be a regular ring, J an ideal of A and \(A_{0}\) a subring of A such that the canonical homomorphism \(A_{0} \rightarrow A/J\) let it be bijective.

\begin{enumerate}
\def\labelenumi{\alph{enumi})}
\item
  The ideal J is completely secant, the \(A_{0}\) -module \(J/J^{2}\) is projective of finite type, and the ring \(A_{0}\) is regular.
\item
  Let \(\varphi: J/J^{2} \to J\) a section \(A_{0}\) -linear of the canonical surjection \(J \to J/J^{2}\) . The homomorphism of \(A_{0}\) -algebras \(\mathsf{S}_{A_{0}}(J/J^{2}) \longrightarrow A\) extending \(\varphi\) extends to an isomorphism of the completion of the graded ring \(\mathsf{S}_{A_{0}}(J/J^{2})\) on the separated completion of \(A\) for the topology \(J\) -adic.
\end{enumerate}

\begin{enumerate}
\def\labelenumi{\alph{enumi})}
\item
  Let \(\mathfrak{p}\) a prime ideal of A containing J. We have \(\mathfrak{p} = (\mathfrak{p} \cap A_0) \oplus J\) and consequently \(\mathfrak{p}^2 = (\mathfrak{p} \cap A_0)^2 \oplus \mathfrak{p}J\) , hence \(\mathfrak{p}^2 \cap J = \mathfrak{p}J\) . Let \(i\) the canonical injection of \(J\) in \(\mathfrak{p}\) . By what precedes, the map \(i \otimes 1_{A/p}: J \otimes_A A/p \longrightarrow \mathfrak{p} \otimes_A A/p\) is injective, and the same is true of \(i_{\mathfrak{p}} \otimes 1_{\kappa(\mathfrak{p})}: J_{\mathfrak{p}} \otimes_{A_{\mathfrak{p}}} \kappa(\mathfrak{p}) \longrightarrow \mathfrak{p}A_{\mathfrak{p}} \otimes_{A_{\mathfrak{p}}} \kappa(\mathfrak{p})\) . Nakayama's lemma implies that the ideal \(J_{\mathfrak{p}}\) of \(A_{\mathfrak{p}}\) is generated by a subset of a system of coordinates of the regular local ring \(A_{\mathfrak{p}}\) . According to prop. 2 of VIII, § 5, no. 3, the ring \(A_{\mathfrak{p}}/J_{\mathfrak{p}}\) is regular and the ideal \(J_{\mathfrak{p}}\) is completely secant. Thus \(J\) is completely secant, the ring \(A_0\) , isomorphic to \(A/J\) is regular, and the \(A_0\) -module \(J/J^2\) is projective of finite type by th. 1 (no. 2).
\item
  It follows from Thm. 1 that the canonical homomorphism \(\beta_{\mathrm{J}}: \mathsf{S}_{\mathrm{A}_{0}}(\mathrm{J}/\mathrm{J}^{2}) \longrightarrow \mathrm{gr}_{\mathrm{J}}(\mathrm{A})\) is bijective. Let \(f: \mathsf{S}_{\mathrm{A}_{0}}(\mathrm{J}/\mathrm{J}^{2}) \longrightarrow \mathrm{A}\) the homomorphism of \(\mathrm{A}_{0}\) -algebras extending the map \(\mathrm{A}_{0}\) -linear \(\varphi: \mathrm{J}/\mathrm{J}^{2} \longrightarrow \mathrm{J}\) . If A is endowed with the J-adic filtration and \(\mathsf{S}_{\mathrm{A}_{0}}(\mathrm{J}/\mathrm{J}^{2})\) with the filtration associated with its grading, \(\beta_{\mathrm{J}}\) is identified with the homomorphism deduced from \(f\) by passing to the associated graded objects. Assertion b) then follows from III, § 2, no. 8, cor. 3 of th. 1.
\end{enumerate}

